%══════════════════════════════════════════════════════════════════════════════
% STORY-CONTENT.TEX — Historia #001: Score Matching y Difusión Estocástica
%
% Único archivo a editar por historia.
% \providecommand (en lugar de \newcommand) evita errores si el archivo
% se carga más de una vez accidentalmente.
%══════════════════════════════════════════════════════════════════════════════

% ─── Identidad ───────────────────────────────────────────────────────────────
\providecommand{\StoryAccentColor}{CtpMauve}
\providecommand{\StoryNumber}{001}
\providecommand{\StoryDate}{2025}
\providecommand{\StorySeriesTag}{QbitLab Stories}
\providecommand{\StoryCategory}{Modelos Generativos}

% ─── Textos ──────────────────────────────────────────────────────────────────
\providecommand{\StoryTitle}{%
  \textquestiondown C\'OMO APRENDE\\ UNA RED A\\ GENERAR RUIDO?%
}
\providecommand{\StorySubtitle}{\textbf{Score matching y difusi\'on estoc\'astica}}
\providecommand{\StoryModelLabel}{Objetivo de entrenamiento (denoising score matching)}

\providecommand{\StoryAbstract}{%
  Los modelos de difusi\'on aprenden a invertir un proceso de corrupci\'on
  gaussiana. El truco central: en lugar de estimar $p(x)$ directamente
  ---intractable--- se estima su gradiente logar\'itmico, el \emph{score}
  $\nabla_x \log p(x)$. Gracias a la identidad de Vincent (2011), esto
  equivale a predecir el ruido a\~nadido: regresi\'on pura.%
}

% ─── Ecuación héroe ──────────────────────────────────────────────────────────
% Renderizada a ~16pt en display mode. Usa macros de equation-styles.tex.
\providecommand{\StoryEquation}{%
  \eqfn{\mathcal{L}}(\equ{\theta})
  \;=\;
  \eqcon{\tfrac{1}{2}}\,
  \eqop{\mathbb{E}}_{\!\eqk{x}}
  \!\left[
    \bigl\|
      \eqfn{s_{\equ{\theta}}}(\eqk{x})
      -
      \eqop{\nabla_{\!\eqk{x}}} \log \eqk{p(x)}
    \bigr\|^2
  \right]%
}

% ─── Figura ───────────────────────────────────────────────────────────────────
\providecommand{\StoryFigure}{%
  \begin{tikzpicture}
    \begin{axis}[
      catppuccin-clean,
      width       = 10cm,
      height      = 6cm,
      xlabel      = {$t$},
      ylabel      = {$\sigma(t)$},
      title       = {Programaci\'on de ruido},
      domain      = 0:1,
      samples     = 80,
      no markers,
      legend pos  = north west,
      title style      = {font=\fontsize{7}{8}\selectfont\bfseries,
                          color=InkDark},
      label style      = {font=\fontsize{6.5}{7}\selectfont,
                          color=InkMid},
      tick label style = {font=\fontsize{6}{7}\selectfont,
                          color=InkLight},
      legend style     = {font=\fontsize{6}{7}\selectfont,
                          draw=InkFaint, fill=white, text=InkMid},
      axis line style  = {color=InkFaint},
      tick style       = {color=InkFaint},
    ]
      \addplot+[thick] {sqrt(1 - exp(-5*x))};
      \addlegendentry{VP-SDE}
      \addplot+[thick] {x};
      \addlegendentry{VE-SDE}
      \addplot+[thick, dashed] {sqrt(x*(1-x))};
      \addlegendentry{Sub-VP}
    \end{axis}
  \end{tikzpicture}%
}

\providecommand{\StoryFigCaption}{%
  Tres esquemas de programaci\'on de ruido. El score $s_\theta(x,t)$
  se entrena para cada nivel $t \in [0,1]$.%
}

% ─── Body ─────────────────────────────────────────────────────────────────────
\providecommand{\StoryBody}{%
  Durante la generaci\'on, el score guia una SDE de Langevin en tiempo
  revertido: el proceso de Ornstein--Uhlenbeck inverso. La curva VP-SDE
  preserva la varianza total; VE-SDE la hace crecer. Ambas convergen a
  la misma distribuci\'on cuando el score es exacto.%
}

% ─── Referencias ─────────────────────────────────────────────────────────────
\providecommand{\StoryRefs}{%
  Song et al.\ (2021) \textit{Score-Based Generative Modeling through SDEs.}
  ICLR.\\
  Vincent (2011) \textit{A Connection Between Score Matching
  and Denoising Autoencoders.} Neural Comput.%
}